\chapter{Mô hình hóa và giải pháp đề xuất}

\section{Tổng quan kiến trúc}

\noindent\textit{Nguồn tham chiếu:}
\url{https://arxiv.org/abs/2506.06147} \textbar{}
\url{https://dl.acm.org/doi/10.14778/3659437.3659454} \textbar{}
\url{https://dl.acm.org/doi/10.14778/3717755.3717761}

WAVES được đề xuất như kiến trúc stream-first với chiến lược trạng thái lai. Kiến trúc tổng thể gồm các thành phần chính:

\begin{enumerate}
  \item \textbf{StreamIngestion}: tiếp nhận dữ liệu từ CSV/JSON/Kafka, gán event-time và ingestion-time, chuẩn hóa schema.
  \item \textbf{WindowManager}: gán bản ghi vào sliding window, tính delta insert/delete khi cửa sổ trượt.
  \item \textbf{BasicDQChecks}: completeness, validity, range -- kiểm tra inline trên luồng chính, $O(1)$ per event.
  \item \textbf{LogicalEngine}: duy trì ngữ cảnh thống kê bằng EMA, sinh elastic bounding boxes tự thích nghi với concept drift.
  \item \textbf{SharedRuleOptimizer}: parse DC, gom nhóm luật có equality columns chung, infinite padding, sinh ActiveBox metadata.
  \item \textbf{Rapidash (Batched Traversal)}: KD-tree bulk load per pane, duyệt nhiều box cùng lúc, box dropping prune nhánh không giao.
  \item \textbf{DecisionLayer}: AlertStateStore, provisional/final/retraction alerts, TombstoneManager.
  \item \textbf{Weever (Pane-based Forest)}: pane close = bulk load, O(1) DROP pane khi expire.
  \item \textbf{WatermarkClock + AlertOutput}: watermark advance trigger pane close, chuyển provisional $\rightarrow$ final, phát meta-stream.
\end{enumerate}

Luồng dữ liệu chính:
\begin{center}
Event $\rightarrow$ WindowManager $\rightarrow$ BasicDQChecks $\rightarrow$ LogicalEngine $\rightarrow$
Rapidash $\rightarrow$ DecisionLayer $\rightarrow$ Weever/Watermark $\rightarrow$ AlertOutput
\end{center}

Rapidash đọc snapshot forest từ Weever (stateless về index). Decision nhận candidate từ Rapidash kèm watermark $\rightarrow$ cập nhật Weever. TombstoneManager gắn vòng đời với pane\_id.

\begin{figure}[htbp]
  \centering
  \includegraphics[width=0.96\linewidth]{architecture}
  \caption{Kiến trúc tổng thể WAVES: StreamIngestion $\rightarrow$ WindowManager $\rightarrow$ BasicDQChecks + LogicalEngine $\rightarrow$ Rapidash $\rightarrow$ DecisionLayer $\rightarrow$ Weever (Pane-based Forest) + Watermark $\rightarrow$ AlertOutput. Mũi tên nét đứt thể hiện Rapidash đọc snapshot từ Weever.}
  \label{fig:waves-arch}
\end{figure}

\section{Research Questions}

Bốn câu hỏi nghiên cứu cốt lõi, mỗi câu gắn với giả thuyết, metrics và modules liên quan.

{\small
\begin{longtable}{|L{0.06\textwidth}|L{0.28\textwidth}|L{0.24\textwidth}|L{0.20\textwidth}|L{0.12\textwidth}|}
\hline
\textbf{RQ} & \textbf{Câu hỏi} & \textbf{Giả thuyết} & \textbf{Metrics} & \textbf{Modules} \\
\hline
RQ1 & Throughput vượt nested-loop truyền thống? & KD-Tree + pane-based forest cắt tỉa không gian & Throughput, P99 latency, detection latency & Rapidash, Weever, Ingestion, Window \\
\hline
RQ2 & EMA + Tombstone/Retraction giữ F1 khi drift + late data? & Elastic box tự thu hẹp/mở rộng theo context; provisional vs final; retraction sửa false alert & Precision, Recall, F1 (sau retraction) & LogicalEngine, Decision, Tombstone \\
\hline
RQ3 & Mở rộng khi 50--100 luật DC? & Shared indexing + batched multi-rule traversal & Memory footprint, throughput & SharedRuleOptimizer, Rapidash batched \\
\hline
RQ4 & Trade-off siêu tham số? & $\alpha$ EMA, pane size, dimension cap ảnh hưởng khác nhau & P99 / avg latency, F1, RAM & ElasticBox, Weever, Optimizer \\
\hline
\end{longtable}
}

\section{Các hệ thống đối chuẩn}

Phương pháp đánh giá dùng ablation study: tắt từng cơ chế cốt lõi để tạo các biến thể suy thoái của WAVES. Cách này công bằng hơn so với so sánh trực tiếp với framework ngoài.

{\small
\begin{longtable}{|L{0.18\textwidth}|c|c|c|c|L{0.30\textwidth}|}
\hline
\textbf{Hệ thống} & \rotatebox{70}{KD-Tree} & \rotatebox{70}{Pane} & \rotatebox{70}{EMA} & \rotatebox{70}{Retraction} & \textbf{Mục đích} \\
\hline
NL-Stream & \ding{55} & \ding{55} & \ding{55} & \ding{55} & Chứng minh $O(N^2)$ là nút thắt cổ chai \\
\hline
Single-Tree-DaQ & \ding{51} & \ding{55} & \ding{55} & \ding{55} & Chứng minh 1 tree lớn $\rightarrow$ fragmentation/spike latency \\
\hline
Static-Box-DaQ & \ding{51} & \ding{51} & \ding{55} & \ding{51} & Chứng minh concept drift $\rightarrow$ false positives \\
\hline
WAVES-SingleRule & \ding{51} & \ding{51} & \ding{51} & \ding{51} & Tắt shared optimizer: mỗi luật = tree riêng; đo overhead \\
\hline
Buffer-Wait-DaQ & \ding{51} & \ding{51} & \ding{51} & \ding{55} & Chứng minh detection latency cao khi watermark như barrier \\
\hline
\textbf{WAVES-Full} & \ding{51} & \ding{51} & \ding{51} & \ding{51} & Performance ceiling \\
\hline
\end{longtable}
}

\textbf{Ghi chú ký hiệu:} \ding{51} = có \,|\, \ding{55} = không

\section{Các luật Denial Constraints (DC1--DC3)}

Ba luật benchmark được sử dụng để đo khả năng phát hiện vi phạm logic.

{\small
\begin{longtable}{|L{0.10\textwidth}|L{0.32\textwidth}|L{0.26\textwidth}|L{0.22\textwidth}|}
\hline
\textbf{DC} & \textbf{Ý nghĩa nghiệp vụ} & \textbf{Công thức} & \textbf{Modules chính} \\
\hline
DC1 & Hai xe cùng quãng đường, xe ngắn không fare cao bất thường. & $\neg(Puloc_s = Puloc_t \wedge dist_s < dist_t \wedge fare_s > fare_t)$ & Rapidash / KD-Tree \\
\hline
DC2 & Cùng tuyến, duration chênh ngoài biên độ giao thông hiện tại. & $\neg(Puloc_s = Puloc_t \wedge Doloc_s = Doloc_t \wedge |dur_s - dur_t| > \Delta_t)$ & LogicalEngine / ElasticBox \\
\hline
DC3 & Cùng tuyến, tolls chênh bất thường. & $\neg(Puloc_s = Puloc_t \wedge Doloc_s = Doloc_t \wedge tolls_s \gg tolls_t)$ & SharedRuleOptimizer / Rapidash \\
\hline
\end{longtable}
}

\section{Fault Injection Pipeline}

NYC Taxi không có ground truth cho lỗi logic. Toàn bộ chương thực nghiệm dựa trên pipeline tiêm nhiễm có kiểm soát, tạo ra benchmark vừa có nhãn sai phạm rõ ràng, vừa mô phỏng được điều kiện streaming thực tế.

{\small
\begin{longtable}{|L{0.12\textwidth}|L{0.28\textwidth}|L{0.34\textwidth}|L{0.16\textwidth}|}
\hline
\textbf{Bước} & \textbf{Nội dung} & \textbf{Kỹ thuật} & \textbf{Output} \\
\hline
B1 -- Base Prep & Chuẩn hóa schema, loại lỗi vật lý rõ ràng. & Filter null, format, range vượt ngưỡng vật lý. & Tập dữ liệu sạch \\
\hline
B2 -- Drift Injection & Mô phỏng concept drift (kẹt xe, bão tuyết). & Tăng trip\_duration theo time-of-day pattern. & Tập có drift \\
\hline
B3 -- Fraud Injection & Tiêm DC1--DC3 trên tỷ lệ nhỏ bản ghi. & Mutate fare\_amount, trip\_duration, tolls\_amount. & Ground truth positives \\
\hline
B4 -- Late/OoO & 90\% on-time, 10\% late 120--300s. & Sinh ingestion\_time khác event\_time, sort theo ingestion. & Benchmark stream \\
\hline
\end{longtable}
}

Nguồn: Icewafl \cite{edbt2025icewafl} (EDBT 2025, DOI: 10.48786/edbt.2025.64).

\section{Chỉ số đo lường}

Hai nhóm metrics độc lập: hiệu năng hệ thống và chất lượng cảnh báo.

{\small
\begin{longtable}{|L{0.15\textwidth}|L{0.22\textwidth}|L{0.22\textwidth}|L{0.31\textwidth}|}
\hline
\textbf{Nhóm} & \textbf{Chỉ số} & \textbf{Ý nghĩa} & \textbf{Dùng trả lời} \\
\hline
Hiệu năng & Throughput & Số bản ghi xử lý được mỗi giây & RQ1, RQ3 \\
\hline
Hiệu năng & P99 Latency & Latency ở percentile 99 & RQ1, RQ4 \\
\hline
Hiệu năng & Detection Latency & Thời gian từ violation $\rightarrow$ final alert & RQ2 \\
\hline
Hiệu năng & Memory Footprint & RAM phân biệt active/cache/tombstone & RQ3, RQ4 \\
\hline
Cảnh báo & Precision & TP/(TP+FP) sau retraction & RQ2, RQ4 \\
\hline
Cảnh báo & Recall & TP/(TP+FN) & RQ2, RQ4 \\
\hline
Cảnh báo & F1 Score & $2PR/(P+R)$ & RQ2, RQ4 \\
\hline
Cảnh báo & FP Reduction Rate & $(FP_{\text{no\_delay}} - FP_{\text{delay}})/FP_{\text{no\_delay}}$ & RQ2 \\
\hline
\end{longtable}
}

\textbf{Lưu ý chấm điểm trong streaming có retraction:} không chốt Precision/Recall chỉ dựa trên provisional alerts [+1]. Một cảnh báo chỉ được xem là còn tồn tại nếu sau giai đoạn quan sát nó không bị triệt tiêu bởi tín hiệu retraction [--1].

\section{Sensitivity Analysis (E1--E3)}

Ngoài benchmark chính, ba thực nghiệm độ nhạy chỉ ra vùng tham số tối ưu.

{\small
\begin{longtable}{|L{0.10\textwidth}|L{0.22\textwidth}|L{0.18\textwidth}|L{0.18\textwidth}|L{0.22\textwidth}|}
\hline
\textbf{Exp} & \textbf{Tham số} & \textbf{Giữ cố định} & \textbf{Metrics} & \textbf{Mong đợi} \\
\hline
E1 & Pane size: 30s/1m/5m/10m & Window 10 phút & P99, avg latency, throughput & Pane nhỏ: nhiều overhead; pane lớn: tree lớn, box dropping kém. Sweet spot: 1--5 phút. \\
\hline
E2 & $\alpha$ (EMA): 0.01/0.05/0.1/0.2 & Tập có drift & Precision, Recall, F1 & $\alpha$ nhỏ: thích nghi chậm, false positives khi drift; $\alpha$ lớn: over-adapt. Sweet spot: 0.05. \\
\hline
E3 & $k_{\max}$: 2/3/5/10 & 100 luật DC & RAM, throughput & $k_{\max}$ nhỏ: box dropping hiệu quả; $k_{\max}$ lớn: curse of dimensionality. DC1--DC3 chỉ cần $k=2$. \\
\hline
\end{longtable}
}

\chapter{Thực nghiệm và đánh giá chi tiết}

\section{Môi trường và dữ liệu}

Tập dữ liệu nền: NYC Taxi Yellow Cab (CSV/Parquet). Các trường ưu tiên: tpep\_pickup/dropoff\_datetime, PULocationID, DOLocationID, trip\_distance, fare\_amount, total\_amount, trip\_duration, tolls\_amount, RatecodeID.

Bản thiết kế giả định kiến trúc stream processing tiêu chuẩn với Kafka ở lớp ingest. Lõi xử lý có thể hiện thực bằng Python kết hợp engine hiệu năng cao ở lớp thấp hơn để tránh điểm nghẽn do runtime Python.

\section{Các kịch bản kiểm thử}

\begin{itemize}
  \item \textbf{Clean stream}: đo chi phí vận hành cơ sở khi không có vi phạm.
  \item \textbf{Logic error without delay}: đo khả năng phát hiện vi phạm logic khi dữ liệu đến đúng thứ tự.
  \item \textbf{Network delay / late arrivals}: đo khả năng giảm false alerts khi dữ liệu đúng nhưng đến muộn.
\end{itemize}

\section{Kết quả kỳ vọng theo từng RQ}

{\small
\begin{longtable}{|L{0.10\textwidth}|L{0.36\textwidth}|L{0.44\textwidth}|}
\hline
\textbf{RQ} & \textbf{Kết quả kỳ vọng} & \textbf{Dấu hiệu thành công} \\
\hline
RQ1 & WAVES vượt NL-Stream về throughput và giữ P99 ổn định hơn Single-Tree. & Đường throughput cao và phẳng; ít latency spikes. \\
\hline
RQ2 & WAVES giữ Precision và F1 cao hơn Static-Box và Buffer-Wait trong dữ liệu có drift + late data. & False positives giảm mạnh; detection latency thấp. \\
\hline
RQ3 & WAVES mở rộng tốt hơn khi tăng số luật lên 50--100. & RAM tăng chậm; throughput không suy sụp. \\
\hline
RQ4 & Tồn tại vùng tham số tối ưu thay vì một cực trị duy nhất. & Biểu đồ sensitivity thể hiện vùng cân bằng rõ ràng. \\
\hline
\end{longtable}
}

\section{Benchmark Pipeline End-to-End}

Quy trình chạy benchmark gồm bốn pha.

{\small
\begin{longtable}{|L{0.12\textwidth}|L{0.40\textwidth}|L{0.38\textwidth}|}
\hline
\textbf{Pha} & \textbf{Việc cần làm} & \textbf{Artefacts đầu ra} \\
\hline
Pha 1 & Chuẩn bị dữ liệu nền và sinh benchmark stream (B1--B4) & Dataset đã gắn ground truth, event time, ingestion time \\
\hline
Pha 2 & Chạy toàn bộ baselines và WAVES-Full trên cùng cấu hình phần cứng & Log runtime, metrics thô, alert streams \\
\hline
Pha 3 & Tổng hợp kết quả theo RQ & Bảng so sánh, biểu đồ, phần diễn giải \\
\hline
Pha 4 & Chạy sensitivity analysis (E1--E3) & Biểu đồ trade-off và cấu hình khuyến nghị \\
\hline
\end{longtable}
}
